\documentclass[aspectratio=169,11pt]{beamer}
% Shared Beamer style for practice contest solution walkthroughs.
\usepackage[utf8]{inputenc}
\usepackage[T1]{fontenc}
\usepackage{lmodern}
\usepackage{amsmath,amssymb}
\usepackage{xcolor}
\usepackage{hyperref}
\usepackage{listings}

\definecolor{CPNavy}{HTML}{172033}
\definecolor{CPBlue}{HTML}{2563EB}
\definecolor{CPGreen}{HTML}{16A34A}
\definecolor{CPOrange}{HTML}{F97316}
\definecolor{CPGray}{HTML}{64748B}
\definecolor{CPLight}{HTML}{F8FAFC}
\definecolor{CPCodeBg}{HTML}{F1F5F9}
\definecolor{CPCodeRule}{HTML}{CBD5E1}

\usetheme{default}
\usefonttheme{professionalfonts}
\setbeamertemplate{navigation symbols}{}
\setbeamersize{text margin left=0.65cm,text margin right=0.65cm}
\setbeamertemplate{itemize item}{\color{CPBlue}$\blacktriangleright$}
\setbeamertemplate{itemize subitem}{\color{CPGray}$\triangleright$}

\setbeamercolor{normal text}{fg=CPNavy,bg=white}
\setbeamercolor{frametitle}{fg=white,bg=CPNavy}
\setbeamercolor{title}{fg=white,bg=CPNavy}
\setbeamercolor{subtitle}{fg=CPLight,bg=CPNavy}
\hypersetup{colorlinks=true,urlcolor=CPBlue}

\setbeamertemplate{frametitle}{%
  \nointerlineskip%
  \begin{beamercolorbox}[wd=\paperwidth,ht=0.88cm,dp=0.22cm,leftskip=0.55cm,rightskip=0.35cm]{frametitle}
    \usebeamerfont{frametitle}\insertframetitle\hfill\small\insertframenumber/\inserttotalframenumber
  \end{beamercolorbox}%
}

\setbeamertemplate{footline}{%
  \leavevmode%
  \hbox{%
    \begin{beamercolorbox}[wd=.58\paperwidth,ht=2.4ex,dp=1ex,leftskip=.5cm]{author in head/foot}%
      \color{CPGray}\insertshorttitle
    \end{beamercolorbox}%
    \begin{beamercolorbox}[wd=.42\paperwidth,ht=2.4ex,dp=1ex,rightskip=.5cm plus1fil]{date in head/foot}%
      \hfill\color{CPGray}\insertshortauthor
    \end{beamercolorbox}%
  }%
}

\setbeamertemplate{title page}{%
  \vspace*{1.25cm}
  \begin{beamercolorbox}[wd=\paperwidth,sep=0.65cm,leftskip=0.15cm,rightskip=0.15cm]{title}
    {\color{white}\Huge\bfseries\inserttitle\par}
    \vspace{0.35cm}
    {\color{CPLight}\Large\insertsubtitle\par}
    \vspace{0.6cm}
    {\color{CPLight}\large\insertauthor\par}
  \end{beamercolorbox}
  \vfill
}

\newcommand{\sectiontag}[1]{\textbf{\color{CPBlue}#1}}
\newenvironment{tightitemize}{\begin{itemize}\small\setlength{\itemsep}{0.18cm}}{\end{itemize}}
\lstdefinestyle{cpcode}{
  language=Python,
  basicstyle=\ttfamily\tiny,
  keywordstyle=\color{CPBlue}\bfseries,
  commentstyle=\color{CPGray}\itshape,
  stringstyle=\color{CPGreen},
  backgroundcolor=\color{CPCodeBg},
  frame=single,
  framerule=0.25pt,
  rulecolor=\color{CPCodeRule},
  showstringspaces=false,
  columns=fullflexible,
  keepspaces=true,
  breaklines=true,
  breakatwhitespace=false,
  tabsize=4,
  xleftmargin=0.08cm,
  xrightmargin=0.08cm,
  aboveskip=0.08cm,
  belowskip=0.08cm
}


\title[flagquiz]{flagquiz}
\subtitle{Day 1 solution walkthrough}
\author[Solution Deck]{Competitive Programming Bootcamp}
\date{}

\begin{document}

\begin{frame}[plain]
  \titlepage
\end{frame}

% BEGIN GENERATED STATEMENT INTERPRETATION
\begin{frame}{Interpret the Word Problem}
\small
\sectiontag{Statement excerpts:}
\begin{tightitemize}
  \item \textit{``find the answer that is easiest to change into any other answer''}
  \item \textit{``smallest maximum amount of changes''}
\end{tightitemize}

\vspace{0.18cm}
\sectiontag{Programming interpretation:}
\begin{tightitemize}
  \item Treat each answer as an ordered list of parts split by comma-space.
  \item The distance between two answers is the number of positions where their parts differ.
  \item For each answer, compute its worst distance to all other answers, then print the minimum-worst candidates.
\end{tightitemize}
\end{frame}
% END GENERATED STATEMENT INTERPRETATION







\begin{frame}{Problem Model}
\small
\sectiontag{Textbook connection:} CP4 2.2 d. Array Manipulation, Harder

\vspace{0.25cm}
\sectiontag{Core technique:} 2D arrays and pairwise comparison

\vspace{0.25cm}
\begin{tightitemize}
  \item Each possible answer is a row in a 2D array of strings.
  \item The incongruity of one answer is its maximum difference against any other answer.
  \item The output is every answer with the smallest possible incongruity, preserving input order.
\end{tightitemize}
\end{frame}

\begin{frame}{How to Recognize the Approach}
\begin{tightitemize}
  \item N is small enough for all-pairs comparison.
  \item Comparing two answer rows is just a coordinate-by-coordinate mismatch count.
  \item The problem asks for a minimax value: minimize the worst distance to another answer.
\end{tightitemize}
\end{frame}

\begin{frame}{Algorithm}
\begin{tightitemize}
  \item Read all answer lines and split each line by comma-space.
  \item For each candidate answer i, compare it with every answer j.
  \item Track the largest mismatch count seen for candidate i.
  \item Maintain the global smallest largest-mismatch value.
  \item Print candidates whose largest mismatch equals the global best.
\end{tightitemize}
\end{frame}

\begin{frame}{Walkthrough of the Key State}
\begin{tightitemize}
  \item countDifferences uses zip so corresponding fields are compared directly.
  \item If two rows are identical, their distance is zero.
  \item When a candidate beats the current best, reset the output list.
  \item When a candidate ties the current best, append it to the output list.
\end{tightitemize}
\end{frame}

\begin{frame}{Why the Algorithm Is Correct}
\begin{tightitemize}
  \item Every possible answer is considered as the candidate final answer.
  \item For each candidate, every possible target answer is checked, so the computed maximum is exact.
  \item Choosing the candidates with minimum exact maximum difference matches the required least incongruous alternatives.
\end{tightitemize}
\end{frame}

\begin{frame}{Complexity and Implementation Notes}
\small
\sectiontag{Complexity:} O(N\textasciicircum{}2 * M) time, where M is the number of fields per answer; O(N * M) memory.

\vspace{0.3cm}
\begin{tightitemize}
  \item Keep rows as lists while comparing, then join only when preparing output.
  \item Do not sort the candidates; ties must stay in input order.
  \item The solution is intentionally direct because the constraints make it safe.
\end{tightitemize}
\end{frame}



% BEGIN GENERATED CODE WALKTHROUGH
\begin{frame}[fragile]{Code: Compare Two Answers}
\scriptsize
\sectiontag{Source lines:} \texttt{flagquiz.py:27--35}

\vspace{0.08cm}
\begin{columns}[T,totalwidth=\textwidth]
\begin{column}{0.58\textwidth}
\begin{lstlisting}[style=cpcode]
import sys

def countDifferences(ans1, ans2):
    count = 0
    for op1, op2 in zip(ans1, ans2):
        if op1 != op2:
            count += 1
    return count
\end{lstlisting}
\end{column}
\begin{column}{0.38\textwidth}
\sectiontag{What this chunk does}
\begin{tightitemize}
  \item The helper performs a position-by-position comparison.
  \item zip pairs corresponding fields, matching the statement rule that order matters.
  \item The returned count is the number of edits between two alternatives.
\end{tightitemize}
\end{column}
\end{columns}
\end{frame}

\begin{frame}[fragile]{Code: Parse Alternatives}
\scriptsize
\sectiontag{Source lines:} \texttt{flagquiz.py:38--45}

\vspace{0.08cm}
\begin{columns}[T,totalwidth=\textwidth]
\begin{column}{0.58\textwidth}
\begin{lstlisting}[style=cpcode]
data = sys.stdin.read().split("\n")

N = int(data[1])
output = []
answers = []

for line in data[2 : N + 2]:
    answers.append(line.split(", "))
\end{lstlisting}
\end{column}
\begin{column}{0.38\textwidth}
\sectiontag{What this chunk does}
\begin{tightitemize}
  \item The question line is ignored because only the alternatives affect the answer.
  \item Each answer line becomes a list of parts, so later comparisons are direct list comparisons.
  \item The input guarantees all alternatives have the same number of parts.
\end{tightitemize}
\end{column}
\end{columns}
\end{frame}

\begin{frame}[fragile]{Code: Measure Each Candidate}
\scriptsize
\sectiontag{Source lines:} \texttt{flagquiz.py:47--56}

\vspace{0.08cm}
\begin{columns}[T,totalwidth=\textwidth]
\begin{column}{0.58\textwidth}
\begin{lstlisting}[style=cpcode]
lowest = float("inf")
for i1 in range(len(answers)):
    instanceGreatest = float("-inf")
    for i2 in range(len(answers)):
        if i1 != i2 and answers[i1] == answers[i2]:
            instanceGreatest = 0
        else:
            diffTest = countDifferences(answers[i1], answers[i2])
            if diffTest > instanceGreatest:
                instanceGreatest = diffTest
\end{lstlisting}
\end{column}
\begin{column}{0.38\textwidth}
\sectiontag{What this chunk does}
\begin{tightitemize}
  \item The outer loop selects one answer as the candidate.
  \item The inner loop compares it against every possible answer.
  \item instanceGreatest records the candidate's incongruousity: its worst observed difference.
\end{tightitemize}
\end{column}
\end{columns}
\end{frame}

\begin{frame}[fragile]{Code: Keep the Best Ties}
\scriptsize
\sectiontag{Source lines:} \texttt{flagquiz.py:57--63}

\vspace{0.08cm}
\begin{columns}[T,totalwidth=\textwidth]
\begin{column}{0.58\textwidth}
\begin{lstlisting}[style=cpcode]
    if instanceGreatest < lowest:
        lowest = instanceGreatest
        output = [", ".join(list(map(str, answers[i1])))]
    elif instanceGreatest == lowest:
        output.append(", ".join(list(map(str, answers[i1]))))

sys.stdout.write("\n".join(output))
\end{lstlisting}
\end{column}
\begin{column}{0.38\textwidth}
\sectiontag{What this chunk does}
\begin{tightitemize}
  \item A strictly better candidate resets the answer list.
  \item A tie is appended, preserving input order.
  \item The stored list form is joined back into the required comma-separated output.
\end{tightitemize}
\end{column}
\end{columns}
\end{frame}
% END GENERATED CODE WALKTHROUGH

\begin{frame}{Source}
\small
\sectiontag{Solution file:} \texttt{practice\_contests/day\_1/flagquiz.py}

\vspace{0.3cm}
\sectiontag{Problem:} \url{https://open.kattis.com/problems/flagquiz}

\vspace{0.45cm}
Use this deck with the source code open beside it. The slides explain the reasoning; the code shows the exact input handling and output formatting.
\end{frame}

\end{document}
